\newif\ifteacher\teachertrue
\documentclass[11pt,a4paper]{article}

\usepackage[margin=2.3cm]{geometry}
\usepackage{amsmath,amssymb,amsthm}
\usepackage[T1]{fontenc}
\usepackage{tgschola}
\usepackage{fouriernc}
\usepackage{microtype}
\usepackage{tfrupee}
\linespread{1.04}
\usepackage{xcolor}
\usepackage{titlesec}
\usepackage{tikz}
\usetikzlibrary{angles,quotes,calc,arrows.meta}
\usepackage{enumitem}
\usepackage{array}
\usepackage{needspace}
\usepackage{fancyhdr}
\usepackage[hidelinks]{hyperref}

\definecolor{accent}{RGB}{38,52,70}
\definecolor{panel}{RGB}{236,238,241}

\newtheoremstyle{brand}{9pt}{9pt}{\itshape}{}{\bfseries\color{accent}}{.}{ }{}
\newtheoremstyle{branddef}{9pt}{9pt}{}{}{\bfseries\color{accent}}{.}{ }{}
\theoremstyle{brand}
\newtheorem{theorem}{Theorem}
\newtheorem{corollary}{Corollary}
\theoremstyle{branddef}
\newtheorem{definition}{Definition}
\newtheorem{example}{Example}[section]
\newtheorem{illus}{Illustration}
\newtheorem{remark}{Remark}
\newtheorem*{activity}{Activity}

\newcommand{\tr}{\triangle}
\newenvironment{exercise}[1]
  {\par\medskip\Needspace{7\baselineskip}\noindent{\large\bfseries\color{accent}Exercise #1}
   \par\smallskip
   \begin{enumerate}[leftmargin=*,itemsep=4pt]}
  {\end{enumerate}}
\definecolor{muted}{RGB}{70,72,76}
\newcommand{\nl}{\\[3pt]}
\ifteacher
\newenvironment{answerblock}
  {\par\smallskip\Needspace{10\baselineskip}\begingroup\small\linespread{1.15}\selectfont\color{muted}\noindent\textit{Answers.}\par}
  {\par\endgroup\medskip}
\newcommand{\ans}[1]{\par\smallskip\Needspace{9\baselineskip}\noindent\textbf{#1}\ }
\else
\usepackage{comment}
\excludecomment{answerblock}
\fi
\newcommand{\figph}[1]{\begin{center}\fbox{\parbox{.85\linewidth}{\small FIGURE: #1}}\end{center}}
% right-angle mark at vertex #1 between the directions to #2 and #3
\newcommand{\ramark}[3]{\coordinate (ra1) at ($(#1)!2mm!(#2)$);\coordinate (ra2) at ($(#1)!2mm!(#3)$);\draw[thin] (ra1) -- ($(ra1)+(ra2)-(#1)$) -- (ra2);}
% side label: #4 placed beside segment #2--#3, on the right-hand side when walking from #2 to #3
% (for a triangle named anticlockwise this is the outside); optional #1 = position along the segment
\newcommand{\slab}[4][0.5]{\path (#2) -- (#3) node[pos=#1, auto=right, inner sep=2.5pt] {#4};}
% dimension line for a whole side: drawn #3 away from segment #1--#2 on its right-hand side, labelled #4
\newcommand{\dimlab}[4]{\draw[thin,{Bar[width=5pt]}-{Bar[width=5pt]}] ($(#1)!#3!-90:(#2)$) -- ($(#2)!#3!90:(#1)$) node[midway, auto=right, inner sep=2pt] {#4};}
\newcommand{\exfig}[1]{\begin{center}\begin{tikzpicture}[scale=.75, every node/.style={font=\small}]#1\end{tikzpicture}\end{center}}
% rough answer diagram (teacher copy only; it sits inside answerblock)
\newcommand{\ansfig}[1]{\begin{center}\begin{tikzpicture}[scale=.6, every node/.style={font=\footnotesize}, draw=muted]#1\end{tikzpicture}\end{center}}
\newcommand{\rememberbox}[1]{%
  \par\medskip\noindent
  \colorbox{panel}{%
    \parbox{\dimexpr\linewidth-2\fboxsep\relax}{%
      \medskip
      \textbf{\color{accent}\large Things to Remember}
      \par\smallskip #1 \medskip}}
  \par\medskip}
\ifteacher
\newcommand{\facilitatornote}[1]{%
  \par\medskip\noindent
  \colorbox{panel}{%
    \parbox{\dimexpr\linewidth-2\fboxsep\relax}{%
      \smallskip
      {\small\textbf{Note for the teacher.}\ #1}
      \smallskip}}
  \par\medskip}
\else
\newcommand{\facilitatornote}[1]{}
\fi

\newcommand{\ModuleTitle}{CIRCLES}
\newcommand{\PartTitle}{Lesson 1}
\newcommand{\ModuleDescriptor}{Grade 9 mathematics facilitator copy}

\pagestyle{plain}
\setlength{\parskip}{4pt}
\setlength{\parindent}{0pt}
% keep tall fractions in running text from touching the line above or below
\setlength{\lineskiplimit}{2pt}
\setlength{\lineskip}{3pt}


\titleformat{\section}{\large\bfseries}{\thesection}{0.8em}{}
\titlespacing*{\section}{0pt}{16pt}{6pt}
\titleformat{\subsection}{\normalsize\bfseries}{\thesubsection}{0.6em}{}
% plain ruled panels instead of coloured boxes
\renewcommand{\colorbox}[2]{\par\smallskip\noindent\textcolor{black!35}{\rule{\linewidth}{0.4pt}}\par\nointerlineskip\vspace{3pt}\noindent#2\par\vspace{2pt}\noindent\textcolor{black!35}{\rule{\linewidth}{0.4pt}}\par\smallskip}
\begin{document}
\noindent{\LARGE Circles}\hfill{\large Lesson 1}\par\vspace{3pt}
\noindent\rule{\linewidth}{0.6pt}\par\medskip

\colorbox{panel}{%
  \parbox{\dimexpr\linewidth-2\fboxsep\relax}{%
    \small\textbf{Teacher's copy: overview.}
    \begin{itemize}[leftmargin=1.6em,itemsep=1pt,topsep=2pt]
      \item \textbf{Aims.} Students use the terms of a circle correctly. They explain why equal chords subtend equal angles at the centre, and the converse. They explain why the perpendicular from the centre bisects a chord, and the converse. They find chord lengths and distances with the Pythagoras theorem.
      \item \textbf{Syllabus.} CBSE Class IX 2026--27, Unit IV Geometry (25 of 80 marks). The syllabus asks students to \emph{explain why} each result holds. The explanations here use congruence (SSS, SAS, RHS).
      \item \textbf{Method.} Students discover each result by solving a short problem first. Give 3--4 minutes, then discuss the answers and state the theorem. Students need a ruler and a compass.
      \item \textbf{Timing.} Section 1: 5 min. Section 2: 15 min. Section 3: 20 min. Practice: 15 min. Exit question: 5 min.
      \item \textbf{Next lesson.} The circle through three points and the circumcircle of a triangle.
    \end{itemize}}}

\section{Parts of a circle \normalfont\small(5 min)}

Start with familiar circles: the wheel of a bullock cart and the Ashoka Chakra on the national flag. Ask: what is the same about every spoke of a wheel?

\begin{definition}
A \textbf{circle} is the set of all points in a plane at a fixed distance from a fixed point. The fixed point is the \textbf{centre}. The fixed distance is the \textbf{radius}.
\end{definition}

\begin{itemize}[leftmargin=1.8em,itemsep=1pt]
  \item A \textbf{chord} is a line segment that joins two points on the circle.
  \item A \textbf{diameter} is a chord through the centre. Diameter $=2\times$ radius.
  \item An \textbf{arc} is a part of the circle. A chord cuts the circle into a \textbf{minor arc} and a \textbf{major arc}.
  \item A \textbf{segment} is the region between a chord and an arc.
  \item A \textbf{sector} is the region between two radii and an arc.
\end{itemize}

\exfig{
  \coordinate (O) at (0,0);
  \fill[black!12] (O) -- (20:2.4) arc (20:75:2.4) -- cycle;
  \fill[black!28] (200:2.4) arc (200:290:2.4) -- cycle;
  \draw (O) circle (2.4);
  \draw (200:2.4) -- (290:2.4);
  \draw (O) -- (20:2.4); \draw (O) -- (75:2.4);
  \draw (150:2.4) -- (330:2.4);
  \fill (O) circle (1.5pt);
  \node[below=2pt] at (O) {$O$};
  \node[left] at (200:2.4) {$P$};
  \node[below] at (290:2.4) {$Q$};
  \node[left] at (150:2.4) {$B$};
  \node[right] at (330:2.4) {$C$};
  \node[above right] at (20:2.4) {$A$};
  \node[above] at (75:2.4) {$D$};
  \node[font=\footnotesize] at (47:1.6) {sector};
  \node[font=\footnotesize,align=center] at (245:2.0) {segment};
  \node[right,font=\footnotesize] at (2.9,0.6) {$OA$: radius};
  \node[right,font=\footnotesize] at (2.9,0.1) {$BC$: diameter};
  \node[right,font=\footnotesize] at (2.9,-0.4) {$PQ$: chord};
}

\facilitatornote{Ask students to draw a circle and name each part on their own drawing. Students often confuse a segment with a sector. A sector always has two radii as its sides; a segment has a chord.}

\section{Equal chords and angles at the centre \normalfont\small(15 min)}

\textbf{Solve first.}
\begin{enumerate}[label=\textbf{P\arabic*.},leftmargin=2.2em,itemsep=3pt]
  \item A circle has centre $O$ and radius $5$ cm. Chords $AB$ and $CD$ are each $6$ cm long. Draw a rough diagram and join $OA$, $OB$, $OC$ and $OD$.
  \begin{enumerate}[label=(\alph*),itemsep=1pt]
    \item Write the three pairs of equal sides in $\tr AOB$ and $\tr COD$.
    \item Are the two triangles congruent? Name the rule.
    \item What can you say about $\angle AOB$ and $\angle COD$?
  \end{enumerate}
  \item In another circle with centre $O$ and radius $5$ cm, $\angle POQ=\angle ROS=60^\circ$.
  \begin{enumerate}[label=(\alph*),itemsep=1pt]
    \item What kind of triangle is $\tr POQ$? Give a reason.
    \item Find the lengths of the chords $PQ$ and $RS$.
    \item What do you notice about the two chords?
  \end{enumerate}
\end{enumerate}

\begin{answerblock}
\ans{P1.} (a) $OA=OC=5$ cm and $OB=OD=5$ cm (radii); $AB=CD=6$ cm (given). (b) Yes, $\tr AOB\cong\tr COD$ by SSS. (c) $\angle AOB=\angle COD$ (corresponding parts of congruent triangles).
\ans{P2.} (a) $OP=OQ=5$ cm, so $\angle OPQ=\angle OQP$. They add up to $180^\circ-60^\circ=120^\circ$, so each is $60^\circ$. The triangle is equilateral. (b) $PQ=5$ cm and, in the same way, $RS=5$ cm. (c) The chords are equal.
\end{answerblock}

\facilitatornote{Let students work on P1 and P2 before you state Theorems 1 and 2. P1 leads to Theorem 1 and P2 leads to its converse.}

\begin{theorem}
Equal chords of a circle subtend equal angles at the centre.
\end{theorem}

\exfig{
  \coordinate (O) at (0,0);
  \coordinate (A) at (100:2.2); \coordinate (B) at (170:2.2);
  \coordinate (C) at (240:2.2); \coordinate (D) at (310:2.2);
  \draw (O) circle (2.2);
  \draw (A) -- (B); \draw (C) -- (D);
  \draw (O) -- (A); \draw (O) -- (B); \draw (O) -- (C); \draw (O) -- (D);
  \pic[draw,angle radius=9pt] {angle=A--O--B};
  \pic[draw,angle radius=9pt] {angle=C--O--D};
  \draw[thin] ($(A)!0.5!(B)!0.1!90:(B)$) -- ($(A)!0.5!(B)!0.1!-90:(B)$);
  \draw[thin] ($(C)!0.5!(D)!0.1!90:(D)$) -- ($(C)!0.5!(D)!0.1!-90:(D)$);
  \fill (O) circle (1.3pt);
  \node[right=2pt] at (O) {$O$};
  \node[above] at (A) {$A$}; \node[left] at (B) {$B$};
  \node[below left] at (C) {$C$}; \node[below right] at (D) {$D$};
}

\textbf{Why it works.} In $\tr AOB$ and $\tr COD$:
$OA=OC$ (radii), $OB=OD$ (radii) and $AB=CD$ (given).
So $\tr AOB\cong\tr COD$ (SSS). Hence $\angle AOB=\angle COD$ (corresponding parts of congruent triangles).

\begin{theorem}[Converse]
If two chords of a circle subtend equal angles at the centre, then the chords are equal.
\end{theorem}

\textbf{Why it works.} In $\tr AOB$ and $\tr COD$: $OA=OC$, $\angle AOB=\angle COD$ and $OB=OD$. So $\tr AOB\cong\tr COD$ (SAS). Hence $AB=CD$.

\begin{example}
In a circle with centre $O$, chords $PQ$ and $RS$ are equal and $\angle POQ=80^\circ$. Find $\angle ROS$ and $\angle ORS$.
\end{example}

\exfig{
  \coordinate (O) at (0,0);
  \coordinate (P) at (60:2); \coordinate (Q) at (140:2);
  \coordinate (R) at (220:2); \coordinate (S) at (300:2);
  \draw (O) circle (2);
  \draw (P) -- (Q); \draw (R) -- (S);
  \draw (O) -- (P); \draw (O) -- (Q); \draw (O) -- (R); \draw (O) -- (S);
  \pic[draw,angle radius=9pt,"$80^\circ$",angle eccentricity=2] {angle=P--O--Q};
  \fill (O) circle (1.3pt);
  \node[right=3pt] at (O) {$O$};
  \node[above right] at (P) {$P$}; \node[above left] at (Q) {$Q$};
  \node[below left] at (R) {$R$}; \node[below right] at (S) {$S$};
}

\textbf{Given:} $PQ=RS$ and $\angle POQ=80^\circ$.\\
\textbf{To find:} $\angle ROS$ and $\angle ORS$.
\begin{enumerate}[leftmargin=1.8em,itemsep=1pt]
  \item $\angle ROS=\angle POQ=80^\circ$ (equal chords subtend equal angles at the centre, Theorem 1).
  \item $OR=OS$ (radii), so $\tr ORS$ is isosceles and $\angle ORS=\angle OSR$.
  \item $\angle ORS+\angle OSR=180^\circ-80^\circ=100^\circ$ (angle sum of a triangle), so $\angle ORS=50^\circ$.
\end{enumerate}
\textbf{Answer:} $\angle ROS=80^\circ$ and $\angle ORS=50^\circ$.

\facilitatornote{Stress that the angles are at the \emph{centre}. Ask: ``Would the result hold if the angles were measured at a point on the circle?'' Leave the question open for a later lesson.}
\section{The perpendicular from the centre to a chord \normalfont\small(20 min)}

\textbf{Solve first.}
\begin{enumerate}[label=\textbf{P\arabic*.},leftmargin=2.2em,itemsep=3pt,start=3]
  \item A circle has centre $O$ and radius $41$ cm. $PQ$ is a chord, and $ON\perp PQ$ with $ON=9$ cm. Draw a rough diagram and join $OP$ and $OQ$.
  \begin{enumerate}[label=(\alph*),itemsep=1pt]
    \item Name the two right triangles. Are they congruent? Name the rule.
    \item Find $PN$ and $NQ$.
    \item What do you notice about the point $N$?
  \end{enumerate}
  \item A circle has centre $O$ and radius $26$ cm. Chord $AB$ is $48$ cm long, and $M$ is the midpoint of $AB$. Join $OA$, $OB$ and $OM$.
  \begin{enumerate}[label=(\alph*),itemsep=1pt]
    \item Show that $\tr OMA\cong\tr OMB$. Name the rule.
    \item $\angle OMA+\angle OMB=\,?$ Use this to find $\angle OMA$.
    \item Find the length $OM$.
  \end{enumerate}
\end{enumerate}

\begin{answerblock}
\ans{P3.} (a) $\tr ONP$ and $\tr ONQ$; $\angle ONP=\angle ONQ=90^\circ$, $OP=OQ$ (radii), $ON$ common, so they are congruent by RHS. (b) $PN^2=41^2-9^2=1681-81=1600$, so $PN=40$ cm, and $NQ=40$ cm. (c) $N$ is the midpoint of $PQ$.
\ans{P4.} (a) $AM=BM$ (given), $OA=OB$ (radii), $OM$ common, so $\tr OMA\cong\tr OMB$ by SSS. (b) $180^\circ$ (linear pair). The two angles are equal, so $\angle OMA=90^\circ$. (c) $AM=24$ cm, $OM^2=26^2-24^2=676-576=100$, so $OM=10$ cm.
\end{answerblock}

\facilitatornote{P3 leads to Theorem 3 and P4 leads to Theorem 4. Part (c) of each problem previews the calculation in the examples below.}

\begin{theorem}
The perpendicular from the centre of a circle to a chord bisects the chord.
\end{theorem}

\exfig{
  \coordinate (O) at (0,0);
  \coordinate (A) at (215:2.2); \coordinate (B) at (325:2.2);
  \coordinate (M) at ($(A)!0.5!(B)$);
  \draw (O) circle (2.2);
  \draw (A) -- (B); \draw (O) -- (M); \draw (O) -- (A); \draw (O) -- (B);
  \ramark{M}{O}{B}
  \draw[thin] ($(A)!0.5!(M)+(0,0.1)$) -- ($(A)!0.5!(M)-(0,0.1)$);
  \draw[thin] ($(M)!0.5!(B)+(0,0.1)$) -- ($(M)!0.5!(B)-(0,0.1)$);
  \fill (O) circle (1.3pt);
  \node[above] at (O) {$O$};
  \node[below left] at (A) {$A$}; \node[below right] at (B) {$B$};
  \node[below] at (M) {$M$};
}

\textbf{Why it works.} Let $OM\perp AB$. In $\tr OMA$ and $\tr OMB$:
$\angle OMA=\angle OMB=90^\circ$, $OA=OB$ (radii) and $OM$ is common.
So $\tr OMA\cong\tr OMB$ (RHS). Hence $AM=BM$.

\begin{theorem}[Converse]
The line joining the centre of a circle to the midpoint of a chord is perpendicular to the chord.
\end{theorem}

\textbf{Why it works.} Let $M$ be the midpoint of $AB$. In $\tr OMA$ and $\tr OMB$: $AM=BM$, $OA=OB$ (radii) and $OM$ is common. So $\tr OMA\cong\tr OMB$ (SSS) and $\angle OMA=\angle OMB$. These two angles form a linear pair, so each is $90^\circ$.

\rememberbox{The \textbf{distance of a chord from the centre} is the length of the perpendicular from the centre to the chord. If the radius is $r$, the chord has length $c$ and its distance from the centre is $d$, then the right triangle $OMA$ gives
\[\left(\dfrac{c}{2}\right)^2+d^2=r^2.\]}

\begin{example}
A chord of a circle of radius $10$ cm is $6$ cm from the centre. Find the length of the chord.
\end{example}

\exfig{
  \coordinate (O) at (0,0);
  \coordinate (A) at (-1.6,-1.2); \coordinate (B) at (1.6,-1.2); \coordinate (M) at (0,-1.2);
  \draw (O) circle (2);
  \draw (A) -- (B); \draw (O) -- (M); \draw (O) -- (A);
  \ramark{M}{O}{B}
  \fill (O) circle (1.3pt);
  \node[above] at (O) {$O$};
  \node[left] at (A) {$A$}; \node[right] at (B) {$B$}; \node[below] at (M) {$M$};
  \node[right] at ($(O)!0.5!(M)$) {$6$};
  \node[above left] at ($(O)!0.5!(A)$) {$10$};
}

\textbf{Given:} radius $OA=10$ cm, $OM\perp AB$, $OM=6$ cm.\\
\textbf{To find:} the length of chord $AB$.
\begin{enumerate}[leftmargin=1.8em,itemsep=1pt]
  \item $AM=MB$ (the perpendicular from the centre bisects the chord, Theorem 3).
  \item $\tr OMA$ is right-angled at $M$, so $AM^2=OA^2-OM^2=100-36=64$ (Pythagoras theorem). So $AM=8$ cm.
  \item $AB=2\times AM=2\times 8=16$ cm.
\end{enumerate}
\textbf{Answer:} the chord is $16$ cm long.

\begin{example}
A chord of length $24$ cm is $9$ cm from the centre of a circle. Find the radius of the circle.
\end{example}

\exfig{
  \coordinate (O) at (0,0);
  \coordinate (A) at (-1.6,-1.2); \coordinate (B) at (1.6,-1.2); \coordinate (M) at (0,-1.2);
  \draw (O) circle (2);
  \draw (A) -- (B); \draw (O) -- (M); \draw (O) -- (A);
  \ramark{M}{O}{B}
  \fill (O) circle (1.3pt);
  \node[above] at (O) {$O$};
  \node[left] at (A) {$A$}; \node[right] at (B) {$B$}; \node[below] at (M) {$M$};
  \node[right] at ($(O)!0.5!(M)$) {$9$};
  \node[above left] at ($(O)!0.5!(A)$) {$r$};
  \node[above] at ($(M)!0.5!(B)$) {$12$};
}

\textbf{Given:} $AB=24$ cm, $OM\perp AB$, $OM=9$ cm.\\
\textbf{To find:} the radius $OA$.
\begin{enumerate}[leftmargin=1.8em,itemsep=1pt]
  \item $AM=\dfrac{24}{2}=12$ cm (Theorem 3).
  \item $OA^2=AM^2+OM^2=144+81=225$ (Pythagoras theorem in $\tr OMA$). So $OA=15$ cm.
\end{enumerate}
\textbf{Answer:} the radius is $15$ cm.

\facilitatornote{Students often use the whole chord in place of half of it. Ask them to mark the midpoint $M$ and the right angle before they calculate.}
\section{Practice \normalfont\small(15 min)}

Work in your exercise book. Draw a rough diagram for every question.

\begin{exercise}{1A}
\item The radius of a circle is 13 cm. A chord is 5 cm from the centre. Find the length of the chord. Draw a rough diagram first.
\item A chord of length 30 cm is drawn in a circle of radius 17 cm. Find the distance of the chord from the centre.
\item Chord $AB$ of a circle with centre $O$ is 48 cm long. $OM \perp AB$ and $OM = 7$ cm. Ravi writes: radius $=\sqrt{48^2+7^2}$. Find his mistake, and find the correct radius.
\item In a circle with centre $O$, chords $AB$ and $CD$ are equal and $\angle AOB = 70^\circ$. Find $\angle COD$ and $\angle OCD$. Give a reason for each step.
\item Draw a circle of radius 5 cm and a chord of length 8 cm in it. Use a set square to draw the perpendicular from the centre to the chord. Measure its length. Then find the length by calculation and compare.
\item A round dining table top has radius 50 cm. A carpenter fixes a straight wooden strip on the underside along a chord, 30 cm from the centre. Find the length of the strip.
\end{exercise}

\facilitatornote{If time is short, set Questions 5 and 6 as homework. Questions 1 to 3 give the main practice in the relation $\left(\dfrac{c}{2}\right)^2 + d^2 = r^2$.}

\begin{answerblock}
\ans{1.} $OM \perp AB$, so $M$ is the midpoint of $AB$ (Theorem 3). In right triangle $OMA$,
\[ AM^2 = OA^2 - OM^2 = 13^2 - 5^2 = 169 - 25 = 144, \]
so $AM = 12$ cm, and $AB = 2 \times 12 = 24$ cm. \textbf{Answer:} 24 cm.
\ansfig{
  \coordinate (O) at (0,0);
  \draw (O) circle (2);
  \coordinate (A) at (-1.85,-0.77);
  \coordinate (B) at (1.85,-0.77);
  \coordinate (M) at (0,-0.77);
  \draw (A) -- (B);
  \draw (O) -- (M);
  \draw (O) -- (A);
  \draw (O) -- (B);
  \ramark{M}{O}{A};
  \node[above] at (O) {$O$};
  \node[below left] at (A) {$A$};
  \node[below right] at (B) {$B$};
  \node[below] at (M) {$M$};
  \node[right] at ($(M)!0.5!(O)$) {$5$};
  \node[below] at ($(A)!0.5!(M)$) {$12$};
  \node[above left] at ($(A)!0.5!(O)$) {$13$};
}

\ans{2.} $AM = \dfrac{30}{2} = 15$ cm. In right triangle $OMA$,
\[ OM^2 = OA^2 - AM^2 = 17^2 - 15^2 = 289 - 225 = 64, \]
so $OM = 8$ cm. \textbf{Answer:} 8 cm.
\ansfig{
  \coordinate (O) at (0,0);
  \draw (O) circle (2);
  \coordinate (A) at (-1.76,-0.94);
  \coordinate (B) at (1.76,-0.94);
  \coordinate (M) at (0,-0.94);
  \draw (A) -- (B);
  \draw (O) -- (M);
  \draw (O) -- (A);
  \draw (O) -- (B);
  \ramark{M}{O}{A};
  \node[above] at (O) {$O$};
  \node[below left] at (A) {$A$};
  \node[below right] at (B) {$B$};
  \node[below] at (M) {$M$};
  \node[right] at ($(M)!0.5!(O)$) {$8$};
  \node[below] at ($(A)!0.5!(M)$) {$15$};
  \node[above left] at ($(A)!0.5!(O)$) {$17$};
}

\ans{3.} Ravi used the whole chord in place of the half chord. $M$ is the midpoint of $AB$, so
\[ AM = \dfrac{48}{2} = 24 \text{ cm}. \]
In right triangle $OMA$,
\[ OA^2 = AM^2 + OM^2 = 24^2 + 7^2 = 576 + 49 = 625, \]
so $OA = 25$ cm. \textbf{Answer:} the radius is 25 cm.
\ansfig{
  \coordinate (O) at (0,0);
  \draw (O) circle (2);
  \coordinate (A) at (-1.92,-0.56);
  \coordinate (B) at (1.92,-0.56);
  \coordinate (M) at (0,-0.56);
  \draw (A) -- (B);
  \draw (O) -- (M);
  \draw (O) -- (A);
  \draw (O) -- (B);
  \ramark{M}{O}{A};
  \node[above] at (O) {$O$};
  \node[below left] at (A) {$A$};
  \node[below right] at (B) {$B$};
  \node[below] at (M) {$M$};
  \node[right] at ($(M)!0.5!(O)$) {$7$};
  \node[below] at ($(A)!0.5!(M)$) {$24$};
  \node[above left] at ($(A)!0.5!(O)$) {$25$};
}

\ans{4.} Equal chords subtend equal angles at the centre, so $\angle COD = \angle AOB = 70^\circ$ (Theorem 1). In triangle $OCD$, $OC = OD$ because both are radii, so the triangle is isosceles and $\angle OCD = \angle ODC$. Therefore
\[ \angle OCD = \dfrac{180^\circ - 70^\circ}{2} = \dfrac{110^\circ}{2} = 55^\circ. \]
\textbf{Answer:} $\angle COD = 70^\circ$, $\angle OCD = 55^\circ$.
\ansfig{
  \coordinate (O) at (0,0);
  \draw (O) circle (2);
  \coordinate (A) at (100:2);
  \coordinate (B) at (170:2);
  \coordinate (C) at (235:2);
  \coordinate (D) at (305:2);
  \draw (A) -- (B);
  \draw (C) -- (D);
  \draw (O) -- (C);
  \draw (O) -- (D);
  \coordinate (mab) at ($(A)!0.5!(B)$);
  \coordinate (mcd) at ($(C)!0.5!(D)$);
  \draw[thin] ($(mab)!0.08!90:(A)$) -- ($(mab)!0.08!-90:(A)$);
  \draw[thin] ($(mcd)!0.08!90:(C)$) -- ($(mcd)!0.08!-90:(C)$);
  \draw (235:0.4) arc (235:305:0.4);
  \draw ($(C)+(0:0.35)$) arc (0:55:0.35);
  \node[above] at (A) {$A$};
  \node[left] at (B) {$B$};
  \node[below left] at (C) {$C$};
  \node[below right] at (D) {$D$};
  \node[above right] at (O) {$O$};
  \node[below] at (270:0.62) {$70^\circ$};
  \node[font=\scriptsize] at ($(C)+(25:0.8)$) {$55^\circ$};
}

\ans{5.} The half chord is $AM = 4$ cm. In right triangle $OMA$,
\[ OM^2 = OA^2 - AM^2 = 5^2 - 4^2 = 25 - 16 = 9, \]
so $OM = 3$ cm. The measured length should be close to 3 cm. \textbf{Answer:} 3 cm.
\ansfig{
  \coordinate (O) at (0,0);
  \draw (O) circle (2);
  \coordinate (A) at (-1.6,-1.2);
  \coordinate (B) at (1.6,-1.2);
  \coordinate (M) at (0,-1.2);
  \draw (A) -- (B);
  \draw (O) -- (M);
  \draw (O) -- (A);
  \draw (O) -- (B);
  \ramark{M}{O}{A};
  \node[above] at (O) {$O$};
  \node[below left] at (A) {$A$};
  \node[below right] at (B) {$B$};
  \node[below] at (M) {$M$};
  \node[right] at ($(M)!0.5!(O)$) {$3$};
  \node[below] at ($(A)!0.5!(M)$) {$4$};
  \node[above left] at ($(A)!0.5!(O)$) {$5$};
}

\ans{6.} The strip is a chord. The half chord is
\[ \sqrt{50^2 - 30^2} = \sqrt{2500 - 900} = \sqrt{1600} = 40 \text{ cm}, \]
so the strip is $2 \times 40 = 80$ cm long. \textbf{Answer:} 80 cm.
\ansfig{
  \coordinate (O) at (0,0);
  \draw (O) circle (2);
  \coordinate (A) at (-1.6,-1.2);
  \coordinate (B) at (1.6,-1.2);
  \coordinate (M) at (0,-1.2);
  \draw[very thick] (A) -- (B);
  \draw (O) -- (M);
  \draw (O) -- (A);
  \ramark{M}{O}{A};
  \node[above] at (O) {$O$};
  \node[below left] at (A) {$A$};
  \node[below right] at (B) {$B$};
  \node[right] at ($(M)!0.5!(O)$) {$30$};
  \node[above left] at ($(A)!0.5!(O)$) {$50$};
  \node[below] at (0,-1.33) {$80$ cm};
}
\end{answerblock}

\section{Exit question \normalfont\small(5 min)}

Answer on a slip of paper.

\begin{enumerate}[label=(\alph*),leftmargin=*]
\item True or false? ``Every line through the midpoint of a chord passes through the centre.'' Give a reason.
\item A chord of length 10 cm is drawn in a circle of radius 13 cm. How far is the chord from the centre?
\end{enumerate}

\begin{answerblock}
\ans{(a)} False. The line joining the centre to the midpoint of a chord is perpendicular to the chord (Theorem 4). A slanted line through the midpoint does not pass through the centre. Only the perpendicular from the centre bisects the chord (Theorem 3). \textbf{Answer:} false.
\ans{(b)} The half chord is 5 cm. In right triangle $OMA$,
\[ d^2 = 13^2 - 5^2 = 169 - 25 = 144, \]
so $d = 12$ cm. \textbf{Answer:} 12 cm.
\ansfig{
  \coordinate (O) at (0,0);
  \draw (O) circle (2);
  \coordinate (A) at (-1.25,-1.56);
  \coordinate (B) at (1.25,-1.56);
  \coordinate (M) at (0,-1.56);
  \draw (A) -- (B);
  \draw (O) -- (M);
  \draw (O) -- (A);
  \ramark{M}{O}{A};
  \node[above right] at (O) {$O$};
  \node[below left] at (A) {$A$};
  \node[below right] at (B) {$B$};
  \node[right] at ($(M)!0.5!(O)$) {$12$};
  \node[below] at ($(A)!0.5!(M)$) {$5$};
  \node[above left] at ($(A)!0.5!(O)$) {$13$};
}
\end{answerblock}
\end{document}
